% Общие стили оформления.
%%% Минималистичный академический стиль: белый фон, тёмно-серый текст,
%%% синий акцент, без рамок и теней. Никаких стандартных тем beamer.

% --- Цветовая палитра ---
\definecolor{bgWhite}{HTML}{FFFFFF}      % фон
\definecolor{fgDark}{HTML}{1A1A1A}       % основной текст
\definecolor{accentBlue}{HTML}{2563EB}   % акцент (заголовки, структура)
\definecolor{secondaryGray}{HTML}{6B7280}% серый (подписи, нижний колонтитул)
\definecolor{blockFill}{HTML}{F3F4F6}    % подложка блоков

\setbeamercolor{normal text}{fg=fgDark,bg=bgWhite}
\setbeamercolor{title}{fg=fgDark}
\setbeamercolor{subtitle}{fg=secondaryGray}
\setbeamercolor{frametitle}{fg=accentBlue}
\setbeamercolor{structure}{fg=accentBlue}
\setbeamercolor{block title}{fg=accentBlue}
\setbeamercolor{block body}{bg=blockFill}
\setbeamercolor{block title alerted}{fg=accentBlue}
\setbeamercolor{block body alerted}{bg=blockFill}
\setbeamercolor{block title example}{fg=accentBlue}
\setbeamercolor{block body example}{bg=blockFill}
\setbeamercolor{footline}{fg=secondaryGray}
\setbeamercolor{caption}{fg=secondaryGray}
\setbeamercolor{caption name}{fg=secondaryGray}
\setbeamercolor{bibliography item}{fg=fgDark}
\setbeamercolor{bibliography entry author}{fg=fgDark}
\setbeamercolor{bibliography entry title}{fg=fgDark}
\setbeamercolor{bibliography entry location}{fg=fgDark}
\setbeamercolor{bibliography entry note}{fg=fgDark}

% Шрифт без засечек --- гротеск Liberation Sans (гарантированно установлен
% в dockertex-образе; Inter и Source Sans 3 в нём отсутствуют).
\setsansfont{Liberation Sans}
\renewfontfamily\cyrillicfontsf{Liberation Sans}

% Размеры полей слайдов
\setbeamersize{text margin left=1.2cm,%
text margin right=1.2cm}

% выключение кнопок навигации
\beamertemplatenavigationsymbolsempty

% Использовать шрифт с засечками для формул, чтобы формулы не ломались
% https://tex.stackexchange.com/a/34267/104425
\usefonttheme[onlymath]{serif}

% --- Размеры шрифтов ---
\setbeamerfont{frametitle}{size=\Large\bfseries}
\setbeamerfont{title}{size=\LARGE\bfseries}
\setbeamerfont{subtitle}{size=\large}
\setbeamerfont{author}{size=\normalsize}
\setbeamerfont{institute}{size=\small}
\setbeamerfont{date}{size=\normalsize}
\setbeamerfont{caption}{size=\footnotesize}
\setbeamerfont{bibliography item}{size=\small}
\setbeamerfont{bibliography entry author}{size=\small}
\setbeamerfont{bibliography entry title}{size=\small}
\setbeamerfont{bibliography entry location}{size=\small}
\setbeamerfont{bibliography entry note}{size=\small}

% --- Заголовок слайда: простой выровненный по левому краю текст ---
% Без декоративной линии/плашки.
\setbeamertemplate{frametitle}{%
  \nointerlineskip%
  \vskip0.45cm%
  {\usebeamerfont{frametitle}%
    \usebeamercolor[fg]{frametitle}%
  \insertframetitle\par}%
  \vskip0.35cm%
}

% --- Номеровать список статей ---
% https://tex.stackexchange.com/a/419506/104425
\setbeamertemplate{bibliography item}{\insertbiblabel}
% или убрать номера
% \setbeamertemplate{bibliography item}{}

% --- Нижний колонтитул: тонкий прогресс-бар + номер страницы ---
% Ширина бара пропорциональна \insertframenumber/\inserttotalframenumber.
% (appendixnumberbeamer уже загружен, поэтому номера слайдов приложения
% не учитываются.) Замечание: опция кадра [plain, noframenumbering] не
% отключает нижний колонтитул — номер слайда при этом не инкрементируется,
% но прогресс-бар и номер по-прежнему рисуются.
\setbeamertemplate{footline}{%
  \begin{beamercolorbox}[wd=\paperwidth,ht=2pt,dp=0pt]{}%
    {\color{accentBlue}%
      \rule{\dimexpr\paperwidth*\insertframenumber/%
    \inserttotalframenumber\relax}{2pt}}%
    {\color{secondaryGray}%
      \rule{\dimexpr\paperwidth-\paperwidth*\insertframenumber/%
    \inserttotalframenumber\relax}{2pt}}%
  \end{beamercolorbox}%
  \begin{beamercolorbox}[wd=\paperwidth,ht=2.25ex,dp=1ex,right]{}%
    {\color{secondaryGray}\scriptsize\insertframenumber}%
    \hspace*{0.8em}%
  \end{beamercolorbox}%
}

% --- Блоки: без рамок, только лёгкая подложка и акцентный заголовок ---
\setbeamertemplate{block begin}{%
  \vskip\medskipamount%
  \begin{beamercolorbox}[wd=\textwidth,sep=0.7em,left]{block body}%
    {\usebeamerfont{block title}\usebeamercolor[fg]{block title}%
    \insertblocktitle\par\vskip0.35em}%
  }
  \setbeamertemplate{block end}{%
  \end{beamercolorbox}%
  \vskip\medskipamount%
}
\setbeamertemplate{block alerted begin}{%
  \vskip\medskipamount%
  \begin{beamercolorbox}[wd=\textwidth,sep=0.7em,left]{block body alerted}%
    {\usebeamerfont{block title alerted}\usebeamercolor[fg]{block
      title alerted}%
    \insertblocktitle\par\vskip0.35em}%
  }
  \setbeamertemplate{block alerted end}{%
  \end{beamercolorbox}%
  \vskip\medskipamount%
}
\setbeamertemplate{block example begin}{%
  \vskip\medskipamount%
  \begin{beamercolorbox}[wd=\textwidth,sep=0.7em,left]{block body example}%
    {\usebeamerfont{block title example}\usebeamercolor[fg]{block
      title example}%
    \insertblocktitle\par\vskip0.35em}%
  }
  \setbeamertemplate{block example end}{%
  \end{beamercolorbox}%
  \vskip\medskipamount%
}

% --- Списки: компактные, акцентные маркеры ---
\setbeamertemplate{itemize item}{\color{accentBlue}\rule{0.9mm}{0.9mm}}
\setbeamertemplate{itemize subitem}{\color{accentBlue}\rule{0.6mm}{0.6mm}}
\setbeamertemplate{itemize subsubitem}{\color{accentBlue}\rule{0.45mm}{0.45mm}}
\setbeamertemplate{enumerate item}{{\color{accentBlue}\insertenumlabel.}}
\setbeamertemplate{enumerate subitem}{{\color{accentBlue}\insertsubenumlabel.}}
\setbeamertemplate{enumerate
subsubitem}{{\color{accentBlue}\insertsubsubenumlabel.}}
\setbeamertemplate{itemize/enumerate body begin}{%
  \setlength\itemsep{0.35em}%
}
\setbeamertemplate{itemize/enumerate subbody begin}{%
  \setlength\itemsep{0.35em}%
}

% --- Подписи к рисункам: маленькие серые, без рамок ---
% Текст подписи не меняем, задаём только цвет и размер.

% https://tex.stackexchange.com/a/291545/104425
\makeatletter
\def\beamer@framenotesbegin{% at beginning of slide
  \usebeamercolor[fg]{normal text}
  \gdef\beamer@noteitems{}%
  \gdef\beamer@notes{}%
}
\makeatother

% вывод на экран заметок к презентации
\ifnumequal{\value{presnotes}}{0}{}{
  \setbeameroption{show notes}
  \ifnumequal{\value{presnotes}}{2}{
    \setbeameroption{show notes on second screen=\presposition}
  }{}
}
